% Source: https://artofproblemsolving.com/community/c3h2140245_point_mania
% Stable identity: aops-topic-2140245-q01
\textbf{(a)} The triangle is isosceles, so $AC=AB=3$ or $AC=BC=7$. Sides $3,3,7$ violate the triangle inequality ($3+3<7$), so $AC=7$.

\textbf{(b)} The triangle has $AC=BC=7$ and base $AB=3$. The altitude from $C$ meets $AB$ at its midpoint, so its length is $\sqrt{49-\dfrac{9}{4}}=\dfrac{\sqrt{187}}{2}$, and
\[[ABC]=\dfrac{1}{2}\cdot3\cdot\dfrac{\sqrt{187}}{2}=\boxed{\dfrac{3\sqrt{187}}{4}}\approx10.26 .\]

\textbf{(c)} For any point $P$ inside the triangle, the triangles $PAB$, $PBC$, $PCA$ fill triangle $ABC$ without overlapping, and their heights from $P$ are the distances $d_{AB}$, $d_{BC}$, $d_{AC}$ from $P$ to the sides.

\begin{asy}
usepackage("amsmath");
size(3.6cm);
defaultpen(fontsize(10pt));
real hh=sqrt(187)/2;
pair A=(-1.5,0), B=(1.5,0), C=(0,hh), P=(0.3,1.8);
pair ft(pair Q, pair X, pair Y){return X+dot(Q-X,Y-X)/dot(Y-X,Y-X)*(Y-X);}
pair F1=ft(P,A,B), F2=ft(P,B,C), F3=ft(P,C,A);
draw(A--B--C--cycle,linewidth(0.9));
draw(P--A,gray); draw(P--B,gray); draw(P--C,gray);
draw(P--F1,blue+dashed); draw(P--F2,blue+dashed); draw(P--F3,blue+dashed);
dot(P); dot(F1,blue); dot(F2,blue); dot(F3,blue);
label("$A$",A,SW); label("$B$",B,SE); label("$C$",C,N); label("$P$",P,dir(60));
label("$3$",(A+B)/2,S); label("$7$",(B+C)/2,E); label("$7$",(A+C)/2,W);
label("\mbox{}",(min(currentpicture,user=true)+max(currentpicture,user=true))/2,black); // keeps the text after the figure black when LaTeX typesets the labels inline
\end{asy}

So
\[[ABC]=\dfrac{1}{2}\left(3\,d_{AB}+7\,d_{BC}+7\,d_{AC}\right).\]
With the stated distances $4$, $2$, $3$ the right side is $\dfrac{1}{2}(12+14+21)=\dfrac{47}{2}=23.5$, but by (b) the area is about $10.26$. This contradiction shows that no such point exists. (More directly: the altitude from $B$ to $AC$ is $\dfrac{2[ABC]}{7}=\dfrac{3\sqrt{187}}{14}\approx2.93$, and no point inside the triangle is farther than that from $AC$, so $d_{AC}=3$ is already impossible.) The answer $\dfrac{47}{2}$ that the area formula in (c) suggests is therefore not meaningful: the configuration in the problem does not exist. $\blacksquare$
